\scanpage{641}
\mspage{29}

\[
 u_*^f\colon M(I)_f\longrightarrow M(I')_f\tag*{(2.9.5)}
\]

\margin{sous-catégories pleines des objets fibrants}

\uline{compatible aux localiseurs} $W(I)_f$ et $W(I')_f$
induits par $W(I),W(I')$ respectivement, donc
induit un foncteur

\[
 u_*^W\colon\operatorname{Ho}_W(I)\unsure{-|\longrightarrow}\operatorname{Ho}_W(I').\tag*{(2.9.6)}
\]

Ce foncteur est un \uline{adjoint à droite} du
foncteur \deleted{\unsure{$u_W^*$}}

\[
 u^*_W\colon\operatorname{Ho}_W(I')\mathrel{-}\operatorname{Ho}_W(I)\tag*{(2.9.7.)}
\]

induit par $u^*$ (2.9.3).

\uline{Démonstration de} 2.9. Pour prouver $1^\circ$),
il faut vérifier \inserted{d'abord} que les conditions préliminaires
(2.8.1 a) b)) sur $M,W$ sont stables
par passage de $M,W$ à $M(I),W(I)$, ce
qui se fait par \unsure{AQT}. \deleted{De plus, \unsure{et de} \unsure{Au théorème 2.6|Au théorème 2.7}
présence des conditions (2.8.2 a) b)), le \unsure{fait}
théorème 2.6 implique alors que (pour \unsure{le|v}}
\deleted{\unsure{Mono}} D'ailleurs on vérifie tout de suite que
\deleted{\unsure{$W(I)$}} $C(I)=\operatorname{Mono}(M(I))$,
et le th. 2.6 ($1^\circ,2^\circ$) implique \inserted{alors} que

\[
\begin{aligned}
 &\deleted{(C(I),C(I)_*}\\
 C(I)&\longleftrightarrow C(I)_*\overset{\mathrm{def}}{=}TF[I]\\
 \underset{\substack{\|\,\mathrm{def}\\ C(I)\cap W(I)}}{TC(I)}
 &\longleftrightarrow(TC(I))_*\overset{\mathrm{def}}{=}F[I]
\end{aligned}
\]

sont des couples de Quillen spéciaux, \unsure{(Étant|étant}
entendu, pour le deuxième, que $TC(I)$ est bien
% archive sheet number: 640
